\documentclass[11pt,a4paper]{article}
\usepackage[margin=2.2cm]{geometry}
\usepackage[T1]{fontenc}
\usepackage{lmodern}
\usepackage{xcolor}
\usepackage{booktabs}
\usepackage{tabularx}
\usepackage{enumitem}
\usepackage{listings}
\usepackage[hidelinks,colorlinks=true,linkcolor=febblue,urlcolor=febblue]{hyperref}
\definecolor{febblue}{HTML}{1C4687}
\definecolor{codebg}{HTML}{F3F4F6}
\lstset{basicstyle=\ttfamily\small, backgroundcolor=\color{codebg}, frame=single, framerule=0pt,
  framesep=4pt, xleftmargin=4pt, breaklines=true, columns=fullflexible, keepspaces=true, showstringspaces=false}
\setlength{\parskip}{4pt}
\setlist{nosep}

\title{\color{febblue}\LARGE FEBAUTO Racing: Competitor Setup on macOS\\[4pt]\large From nothing to a driving car, step by step}
\author{Formula Electric at Berkeley, Autonomous Team}
\date{September 2026}

\begin{document}
\maketitle

\section*{What you are setting up}

A simulator of the RoboRacer 1:10 car runs natively on your Mac. Your code runs in a Docker
container that carries ROS~2 and the bridge to the simulator. You never install ROS on the
Mac. Foxglove shows you the car's sensors live.

\begin{lstlisting}
feb-racing/                 the repo, your clone
  feb-sim                   the launcher
  stack/                    YOUR code (gitignored)
  starter_kit/feb_driver/   the template you copy into stack/
  tracks/                   track folders
~/.feb-sim/app/mac/         the simulator app, downloaded by setup
Docker Desktop              runs the devkit container
\end{lstlisting}

The loop: the container mounts \texttt{stack/} and builds it; the bridge turns the simulator's
stream into ROS~2 topics under \texttt{/autodrive/roboracer\_1/}; your node reads the sensors
and publishes two floats, \texttt{throttle\_command} and \texttt{steering\_command}, both in
$[-1, 1]$; the bridge sends them to the simulator; the car moves.

\section{Install, once}

\begin{enumerate}
  \item \textbf{Docker Desktop.} Install with \texttt{brew install --cask docker-desktop} (or download
  from \url{https://docker.com}), then \texttt{open -a Docker}. Do not install \texttt{docker/tap/sbx};
  that is Docker's agent sandbox tool, not Docker. On first launch it may report ``Failed to install
  Rosetta'': install Rosetta on the Mac first, \texttt{softwareupdate --install-rosetta
  --agree-to-license}, then relaunch Docker Desktop and let it retry. Our images run natively on Apple
  Silicon, so Rosetta is only for other x86 images later. Check with \texttt{docker run --rm hello-world}.
  Leave the app running.
  It is the whale icon in the menu bar. On a Mac, containers live in a small Linux virtual machine
  that exists only while the app is open, so \texttt{docker} commands fail with ``cannot connect to
  the Docker daemon'' if you quit it. Idle it costs almost nothing; tick \emph{Start Docker Desktop
  when you sign in} in its settings. (On Linux, Docker is a background service that starts at boot,
  which is why nothing has to be opened there.)
  \item \textbf{Foxglove.} Download from \url{https://foxglove.dev/download}, open once, sign in with a free account.
  \item \textbf{GitHub CLI.} In Terminal: \texttt{brew install gh}, then \texttt{gh auth login} and follow the prompts.
  \item \textbf{Python 3.} Already on the Mac. If Terminal says otherwise: \texttt{xcode-select --install}.
\end{enumerate}

\section{Clone and set up}

\begin{lstlisting}
git clone https://github.com/Pranman1/feb-racing
cd feb-racing
./feb-sim setup --team "Your Team"
\end{lstlisting}

Pulls the devkit image and downloads the Mac app into \texttt{\textasciitilde/.feb-sim/app/}.
A few minutes. If the app later refuses to open (``damaged'' or ``cannot be opened''), it is
unsigned; run once:

\begin{lstlisting}
xattr -dr com.apple.quarantine ~/.feb-sim/app/mac/"FEB Simulator.app"
codesign --force --deep -s - ~/.feb-sim/app/mac/"FEB Simulator.app"
\end{lstlisting}

\section{Drive manually}

\begin{lstlisting}
./feb-sim run
\end{lstlisting}

Arrow keys drive. Menu top-left: Manual/Autonomous, Track, Cars, Look, camera. Close the
window to stop. Nothing drives by itself without a stack.

\section{Drive autonomously with the starter kit}

\begin{lstlisting}
cp -r starter_kit/feb_driver stack/
./feb-sim run --stack "ros2 launch feb_driver drive.launch.py"
\end{lstlisting}

The container builds the package, the simulator opens, the car drives, Foxglove opens on its
topics.

\textbf{What the starter driver is.} A complete ROS~2 package. The driver is one commented file,
\texttt{stack/feb\_driver/feb\_driver/reactive\_driver.py}, a ``follow the gap'' controller: clean
the lidar scan, blank a safety bubble around the nearest obstacle, aim at the widest gap, pick a
speed from the room ahead, publish steering and throttle. Its numbers are in
\texttt{stack/feb\_driver/config/driver.yaml}.

\textbf{Editing code.} \texttt{stack/} is a plain folder on your Mac; open it in any editor. With
VS~Code: \texttt{code stack/} from the repo folder, or File, Open Folder. Edit, save, run the same
\texttt{./feb-sim run} command again; the container rebuilds your packages in seconds. No plugin
or remote setup: the container reads the files straight from your disk.

\section{Score and post}

\begin{lstlisting}
./feb-sim practice --stack "ros2 launch feb_driver drive.launch.py"
./feb-sim submit
\end{lstlisting}

\texttt{practice} runs one attempt the way a competition does: headless simulator, a warm-up lap
plus ten timed laps, 10~s per collision, 300~s limit; it writes \texttt{runs/<id>/result.json} and a
bag. \texttt{submit} opens a pull request with that file; a check validates and merges it, and your
time appears on the site's practice board as \emph{unverified}. The organiser's nightly rerun of your
registered image on the same track turns it \emph{verified}.

\section{Looking inside: extra terminals and Foxglove}

\subsection*{More terminals into the devkit}

The container keeps running while the simulator is open. Open as many Terminal tabs as you
like and, in each, from the repo folder:

\begin{lstlisting}
./feb-sim shell
\end{lstlisting}

That drops you into a ROS~2 shell inside the container, with the environment already sourced.
Useful commands there:

\begin{lstlisting}
ros2 topic list                                          # everything on the graph
ros2 topic hz /autodrive/roboracer_1/lidar               # is the sensor stream alive, at what rate
ros2 topic echo /autodrive/roboracer_1/throttle_command  # what your node is sending
ros2 topic echo --once /autodrive/roboracer_1/imu
ros2 node list                                           # your node should be here
ros2 node info /driver
ros2 param list /driver                                  # live parameters
ros2 bag record -a -o /home/autodrive_devkit/src/stack/mybag   # lands in your stack/ folder on the Mac
\end{lstlisting}

\texttt{exit} leaves the shell; the container keeps running.

\textbf{Devkit on its own.} The launcher starts the devkit and the simulator together because
that is what you want most of the time, but they are independent. The devkit alone, for
building, replaying bags or running tests without a simulator:
\begin{lstlisting}
docker run --rm -it -p 4567:4567 -v "$PWD/stack:/home/autodrive_devkit/src/stack" \
    ghcr.io/pranman1/feb-devkit bash
\end{lstlisting}
Inside, \texttt{source install/setup.bash} after \texttt{colcon build}. The simulator alone is the
app in \texttt{\textasciitilde/.feb-sim/app/mac}, keyboard driving, no ROS. \texttt{docker exec -it feb-devkit
bash} is the same thing without the launcher. Files you write under
\texttt{/home/autodrive\_devkit/src/stack} inside the container appear in \texttt{stack/} on the
Mac, and the other way round.

\subsection*{Foxglove}

Installed in step~2 of the install list; sign in once. After that, \texttt{./feb-sim run} opens
Foxglove connected to your car a few seconds after the simulator window. If it does not, or
you want a second window: in Foxglove choose \emph{Open connection}, \emph{Foxglove WebSocket},
address \texttt{ws://localhost:8765}.

Panels worth adding (the \emph{+} at the top right of a panel):
\begin{itemize}
  \item \textbf{3D}: add the topic \texttt{/autodrive/roboracer\_1/lidar}; set the display frame to the lidar's frame. The scan appears as points around the car.
  \item \textbf{Image}: \texttt{/autodrive/roboracer\_1/front\_camera}.
  \item \textbf{Plot}: \texttt{throttle\_command.data} and \texttt{steering\_command.data} against time; add \texttt{left\_encoder.position[0]} to watch the wheel.
  \item \textbf{Raw Messages}: any topic, to read exact values.
  \item \textbf{Topics} (left sidebar): every topic with its rate; a dash means nothing is publishing.
\end{itemize}

Save the layout (\emph{Layout}, \emph{Save}) and it comes back next time. Foxglove also opens the
bags that \texttt{practice} records under \texttt{runs/<id>/}, and any bag you record yourself:
\emph{Open local file}.

\section{When something goes wrong}

\begin{lstlisting}
docker logs feb-devkit        # build errors, your node's output
./feb-sim shell               # a ROS 2 terminal inside the container
./feb-sim stop                # kill everything and start over
\end{lstlisting}

Nothing drives: correct without \texttt{--stack}; with one, read \texttt{docker logs}.
Port 4567 in use: \texttt{./feb-sim stop}. Docker pull denied: Docker Desktop is not running.

\section{Keeping up to date}

\begin{lstlisting}
./feb-sim update
\end{lstlisting}

One command pulls the repo (tracks, launcher, docs, new starter packages), pulls a newer devkit
image if there is one, and replaces the simulator app when a new build is released. Your
\texttt{stack/} folder is never touched. When a new starter package appears in
\texttt{starter\_kit/}, such as the cone-track driver, copy it the same way as the first one:

\begin{lstlisting}
cp -r starter_kit/feb_cone_driver stack/
./feb-sim run --track loop_cones --stack "ros2 launch feb_cone_driver drive.launch.py"
\end{lstlisting}

\section{What the car gives and takes}

\begin{tabularx}{\linewidth}{lX}
\toprule
Topic (under \texttt{/autodrive/roboracer\_1/}) & Meaning \\
\midrule
\texttt{lidar} & \texttt{LaserScan}, 1080 beams over 270\textdegree, 0.06 to 10~m \\
\texttt{imu} & orientation, angular velocity, linear acceleration \\
\texttt{left\_encoder}, \texttt{right\_encoder} & wheel angle in rad; radius 0.059~m \\
\texttt{front\_camera} & image \\
\texttt{throttle\_command}, \texttt{steering\_command} & \textbf{you publish these}: \texttt{Float32} in $[-1, 1]$; throttle 0 is a hard brake; steering $\pm 1 = \pm 0.5236$~rad \\
\texttt{ips}, \texttt{odom}, lap and collision counters & ground truth: training only, barred at race time \\
\bottomrule
\end{tabularx}

\smallskip
Publish as many topics of your own as you like; only the two commands reach the car. Use
message stamps or \texttt{use\_sim\_time}, never the wall clock, for speeds and integrals.
Throttle about 0.1 is roughly 2~m/s.

\end{document}
